\begin{explanationbox}
\textbf{公式不是硬记的：先看两段路怎么转。}
把从 $B$ 到 $P$ 的路拆成水平的 $u$ 和竖直的 $v$，两段都转相同的角。
\begin{center}\begin{tikzpicture}[x=1cm,y=1cm,>=Stealth,font=\small]
\path[use as bounding box] (-6.6,-0.6) rectangle (6.6,5.2);

\node at (0,4.8) {先转两段路，再把横、纵变化相加（图取 $u=3,v=2,\theta=30^\circ$）};
\begin{scope}[shift={(-4.8,0.3)}]
\draw[->,gray] (-0.4,0)--(3.7,0) node[right] {$u$方向};
\draw[->,gray] (0,-0.3)--(0,3.2);
\draw[blue!75!black,very thick,->] (0,0)--(3,0) node[midway,below] {$u$};
\draw[orange!85!black,very thick,->] (3,0)--(3,2) node[midway,right] {$v$};
\draw[gray,dashed] (0,0)--(3,2);
\node[below left] at (0,0) {$B$}; \node[above] at (3,2) {$P$};
\draw (2.8,0)--(2.8,0.2)--(3,0.2);
\end{scope}
\begin{scope}[shift={(1.5,0.3)}]
\coordinate (H) at ({3*cos(30)},{3*sin(30)});
\coordinate (Q) at ({3*cos(30)-2*sin(30)},{3*sin(30)+2*cos(30)});
\draw[->,gray] (-0.5,0)--(4,0) node[right] {$x$};
\draw[->,gray] (0,-0.3)--(0,3.9) node[above] {$y$};
\draw[blue!75!black,very thick,->] (0,0)--(H);
\draw[orange!85!black,very thick,->] (H)--(Q);
\draw[teal!75!black,dashed,thick] (0,0)--(Q);
\draw[blue!75!black,densely dotted] (H)--({3*cos(30)},0);
\node[blue!75!black,below] at (1.3,0) {$u\cos\theta$};
\node[blue!75!black,right] at (2.6,0.7) {$u\sin\theta$};
\draw[orange!85!black,densely dotted] (Q)--({3*cos(30)},{3*sin(30)+2*cos(30)})--(H);
\node[orange!85!black,above] at (2.1,3.65) {$-v\sin\theta$};
\node[orange!85!black,right] at (2.6,2.5) {$v\cos\theta$};
\draw[->] (0.65,0) arc(0:30:0.65); \node at (0.95,0.35) {$\theta$};
\fill (H) circle(1.5pt) node[below right] {$H$};
\fill[teal!75!black] (Q) circle(2pt) node[above left] {$Q$};
\node[below left] at (0,0) {$B$};
\draw ($(H)+(-0.173,-0.1)$)--($(H)+(-0.273,0.073)$)--($(H)+(-0.1,0.173)$);
\end{scope}

\end{tikzpicture}\end{center}

蓝色段旋转后的贡献是 $(u\cos\theta,u\sin\theta)$；橙色段向左上方走，贡献是
$(-v\sin\theta,v\cos\theta)$。\textbf{横向相加得 $s$，纵向相加得 $t$：}
\[
(s,t)=(u\cos\theta-v\sin\theta,\ u\sin\theta+v\cos\theta).
\]
图中 $s=3\sqrt3/2-1,\ t=3/2+\sqrt3$，可由两个小直角三角形直接读出。
这里先画 $u,v>0$ 的锐角情形；一般情形用带正负号的坐标表示移动，推导见后文。

\textbf{最后为何加回中心？} $(s,t)$ 是从 $B$ 到 $Q$ 的移动量，因此
$Q=(a+s,b+t)$。减中心和加中心只是在换算点的坐标，真正旋转的对象始终是 $\overrightarrow{BP}$。

\textbf{特殊角的差别只来自正弦、余弦。}
$45^\circ$ 时两者均为 $1/\sqrt2$，出现和、差；$90^\circ$ 时余弦为零，出现交换变号；
$180^\circ$ 时正弦为零、余弦为 $-1$，两个相对坐标都变号。
\end{explanationbox}
